\scanpage{414}
\mspage{153}
\unsure{je voudrais|je voudrait} prouver$^{*}$ que $\bar f$ est une $\unsure{\mathbb D|\mathrm D}$-fibration, propre et lisse. \deleted{\unsure{$^{*}$|Ce}}
% slanted
\margin{\unsure{Le|Ce} résultat fondamental \unsure{de} \unsure{$^{*}$} \unsure{ce que} a été déjà vu, page 170-171}

$\overline{X'}$ est la catégorie des couples
\[
(x',c)\in\mathrm{Ob}\,X'\times\mathrm{Ch}_{\infty}(X)
\]
tels que
\[
\sigma(c)=f(x')
\]
et
\[
\bar f(x',c)=\beta(c)\qquad\text{but de }c.
\]

La fibre de $\deleted{\unsure{f}}\overline{X'},\bar f$ en $x_0\in X$ est donc \deleted{\unsure{formée}}\inserted{la cat.} des couples $(x',c)\in\mathrm{Ob}\,X'\times\mathrm{Ch}_{\infty}(X)$, tels que
\[
\sigma(c)=\unsure{f(x')},\quad\beta(c)=x_0.
\]

\deleted{\unsure{Cette}} On peut regarder encore $\overline{X'}$ comme limite inductive des sous-catégories $\overline{X'}_{n,m}$, \unsure{correspondant} aux couples $(x',c)$ avec $c\in\mathrm{Ch}_{n,m}(X)$. On trouve que pour $\unsure{m}$ \unsure{pair|paire} (direction \unsure{terminale|dominante} de $\tau_{n,m}$ rétrograde) les $\deleted{\unsure{\mathrm{Ch}}}\overline{X'}_{n,m}$ sont \inserted{Cat-}fibrées sur $X$, pour $m$ impair (direction \unsure{terminale|dominante} des $\tau_{n,m}$ directe) \deleted{elles} sont \inserted{Cat-}cofibrées. Donc on conclut encore comme tantôt, \deleted{par} $\overline{X'}$ étant \inserted{(au choix)} limite inductive \deleted{filtrante} de ces catégories lisses sur $X$, ou propres sur $X$,
% archive sheet number: 413
